% Mechanistic Validity — v6 (TMLR restructure skeleton)
% =====================================================================
% v5 (mechanistic_validity_v5.tex) is preserved intact as the SOURCE to port from.
% This v6 is the reorganized shell designed against the 10-reviewer gauntlet:
%   - theory-forward, taxonomy compressed, apparatus moved to supplement_reference.tex
%   - audit reframed as CLAIM-level (not paper-level)
%   - NEW Section 6 "Validating the Framework Itself" (fixes the #1 reject-driver:
%     nothing executed / retrofitting) via preregistration + one ground-truth Validated claim + IRR
% Fill in section by section. \PORT = lift from v5; \TODO = new writing needed.
% =====================================================================
\pdfoutput=1
\documentclass{article}
\usepackage[preprint]{tmlr}

\newcommand{\openreview}{\url{https://openreview.net/forum?id=XXXX}}
\def\month{06}
\def\year{2026}

\usepackage{amsmath,amssymb}
\usepackage{booktabs}
\usepackage{graphicx}
\graphicspath{{figures/}}
\usepackage{xcolor}
\usepackage{hyperref}
\usepackage{cleveref}
\usepackage{enumitem}
\usepackage{placeins}
\usepackage{float}
\usepackage{multirow}
\usepackage{array}

\newcommand{\ie}{i.e.}
\newcommand{\eg}{e.g.}
\newcommand{\cf}{cf.}
\newcommand{\tighttable}{\setlength{\tabcolsep}{4pt}\renewcommand{\arraystretch}{1.05}}

% author-facing skeleton markers — delete before submission.
% Render as clearly-delimited, READABLE notes so a compiled draft shows at a glance
% what is written prose vs. scaffolding (content muted, red/blue brackets around it).
\newcommand{\TODO}[1]{\par\smallskip\noindent\textcolor{red}{\textbf{[\,TODO}}\ \textit{\textcolor{black!65}{#1}}\ \textcolor{red}{\textbf{TODO\,]}}\par\smallskip}
\newcommand{\PORT}[1]{\par\smallskip\noindent\textcolor{blue!65!black}{\textbf{[\,PORT from v5}}\ \textit{\textcolor{black!65}{#1}}\ \textcolor{blue!65!black}{\textbf{PORT\,]}}\par\smallskip}

\usepackage{caption}
\captionsetup[table]{skip=8pt}

\title{Mechanistic Validity: A Theory of Validity for Mechanistic Claims}

\author{\name Elliot Tower \\
      \email elliot@elliottower.ai}

\begin{document}
\maketitle

% =====================================================================
\begin{abstract}
Mechanistic claims are ubiquitous in empirical science---\emph{this circuit implements this computation}, \emph{this gene causes this disease}, \emph{this pathway mediates this outcome}. What makes such a claim valid?

We argue that a mechanistic claim is only as strong as its weakest validity dimension. Five validity types---construct, measurement, internal, external, and interpretive---form a dependency chain: construct validity is a ceiling on all others, so no amount of causal evidence rescues a claim whose target concept is incoherent. This is a theory of validity, not a taxonomy: it predicts that claims cap at the tier corresponding to their weakest dimension, and that single high-scoring metrics cannot compensate for untested dimensions.

We operationalize this theory as Mechanistic Validity, drawing on philosophy of science, measurement theory, and validation methodology from neuroscience, pharmacology, and genetics. We audit sixteen mechanistic claims---thirteen published circuits plus an SAE feature, a steering vector, and Anthropic's J-space ``global workspace'' claim---auditing each \emph{claim} and its evidence base as it now stands, not indicting any single paper, and show the framework discriminates: the IOI circuit reaches Mechanistically Supported but caps there because cross-task specificity remains untested; induction heads reach Triangulated; the gender bias circuit is Disconfirmed on an incoherent construct. We then test the framework against itself: a preregistered prospective audit and a ground-truth toy circuit calibrate its verdicts.
% \TODO{tighten abstract to lead with the 2-sentence contribution delta (R2): unit = causal mechanistic claim (not a benchmark score); description mode sets the evidence bar. Cut the apparatus counts from the abstract.}
\end{abstract}

% =====================================================================
\section{Introduction}
\label{sec:intro}
% GOAL: state the nugget in sentence 1--2, then the DELTA over validity-for-ML work (R2),
% then the two falsifiable predictions, then promise the audit + the self-validation.
\PORT{\S1 opening catalog (induction heads, IOI, greater-than, SAE features) and the
"strong on one dimension, untested on others" paragraph.}
\TODO{Rewrite the opening so the first two sentences are the contribution delta:
(1) the unit of analysis is a causal mechanistic claim, not a benchmark score;
(2) the declared description mode determines which evidence is required.
Then: the weakest-link thesis + two testable predictions. End by promising Sec 4 (audit)
and Sec 6 (self-validation). Keep short; no apparatus dump here.}
\TODO{DELTA / COMPETITOR CHECK (R2): "Causally Grounded Mechanistic Interpretability for LLMs"
(2026) may propose evaluation criteria ADJACENT to the validity tiers --- READ IT and add an
explicit distinction in related work, or it becomes an "already done" reject vector. Bean
(NeurIPS 2025) + MIB (ICML 2026) are already in v5 as capability/method-performance foils; the
causal-mechanistic-claim UNIT stays the delta --- just make sure the 2026 causal-MI paper
doesn't scoop the framing.}

% =====================================================================
\section{The Pattern Across Sciences}
\label{sec:motivation}
% KEEPER — this is the motivation reviewers praised. Port nearly verbatim.
\PORT{Entire \S2 Motivation: dead salmon, cold fusion, BOLD, candidate genes, cardiac
stents, COVID-ML, shortcut learning, emergent abilities, + the MI failure catalog
(\S2.3) and the cross-science table (tab:cross-science) and "The Common Structure".}
\TODO{Trim only for length; keep the cross-field table. End on: MI evaluates methods and
artifacts but has no shared protocol for the CLAIMS built on them --- this paper is that protocol.}
\TODO{ELEVATE the COVID-ML example (DeGrave, Janizek \& Lee 2020): external validation on new
hospital data PASSES while the causal claim ("detects COVID pathology") is FALSE --- shortcuts
don't hurt held-out performance, they just aren't detecting disease. The cleanest single
demonstration that external validity can pass while internal validity fails = the dependency-chain
thesis in one example. Cold-fusion primary: Fleischmann-Pons 1989 (bypassed peer review) ->
Coulomb barrier + missing neutrons/gammas + replication failure (MIT/Caltech/Harwell); Google 2019
still null.}

% =====================================================================
\section{The Framework in Brief}
\label{sec:framework}
% COMPRESSED SPINE. This one section replaces v5's Framework + Methodology + Foundations.
% Contains ONLY what the audit and the "why it matters" sections actually invoke.
% Everything else -> supplement_reference.tex.

\PORT{Dependency-chain equation and the 5 validity-type paragraphs (construct, measurement,
internal, external, interpretive) --- but tightened to ~1 short paragraph each.}
\begin{equation}
\text{Construct} \rightarrow \text{Measurement} \rightarrow \text{Internal} \rightarrow \text{External} \rightarrow \text{Interpretive}
\label{eq:dependency}
\end{equation}
\TODO{One tight paragraph per type. The chain is a CEILING, not a gate sequence.}

\subsection{Description Modes}
\PORT{Compressed description-modes paragraph + tab:modes (the 7-mode table stays --- V1/V4
depend on it). Cut the long prose; keep the table + 3 sentences.}

\subsection{Evidence Families}
\PORT{One paragraph + the 4x2 tab:families ("most published MI lives in one cell").}

\subsection{Criteria and Verdict Tiers}
\PORT{tab:criteria as a COMPACT one-line-each table (grouped by type), and tab:verdicts.
Fix the numbering: I-criteria currently jump I7->I10 (I8, I9 missing) --- renumber or explain.}
\TODO{Add pointer sentence: "Full pass conditions, the 65-metric inventory, protocols, and
synthesis methods are in the Supplementary Reference (also browsable at [website])."}
\TODO{Reconcile all apparatus counts before submission: v5 says 9 vs 4 protocols and
14 vs 15 calibrations in different places. Pick the true numbers.}

\subsection{Severity: Preregistration and Corroboration}
\label{sec:severity}
% Elliot's design decision: exploratory work must not silently reach the top tiers.
% RECOMMENDED FORMULATION (pending final sign-off):
%  - Triangulated and Validated require the KEY criteria to be established by SEVERE tests (Mayo).
%  - Severity is achieved by EITHER (a) preregistration of the prediction, OR (b) risky
%    corroboration: confirmation via a method / distribution / model the claim was NOT built or
%    optimized to pass (independent evidence family, held-out prompts, cross-model).
% REFINED RESOLUTION (the cap is not a NEW penalty --- it FALLS OUT of the existing structure):
%  - Triangulated is DEFINED by convergence across >=3 independent evidence families (C3): it
%    already REQUIRES corroboration. Single-method confirmation-only work already fails E1
%    (>=2 methods) -> caps at Causally Suggestive. So "within-family / confirmation-only evidence
%    caps at Mechanistically Supported" is IMPLIED by the tiers, not bolted on.
%  - Preregistration therefore is NOT a separate tier gate and does NOT let a claim skip
%    cross-family triangulation. It enters two ways: (i) discussed IN PROSE only --- the verdict
%    LABEL stays a clean word ("Triangulated", NOT "Triangulated (preregistered)"; no asterisk ---
%    Elliot's call); each per-claim writeup states in text whether the tier was earned by
%    corroboration and/or preregistration --- and (ii) a severity BOOSTER on individual criterion
%    statuses (Mayo) --- a preregistered specificity test earns "Confirmed" where a post-hoc one
%    earns "Partially confirmed".
%  - Net = Elliot's goal (exploratory work can't silently be Triangulated) WITHOUT a crude gate,
%    fully derived from the framework's own evidence-family + tier structure.
% CAP PROS/CONS (for the writeup + Jensen):
%   moderate cap @ Mechanistically Supported -> principled (Triangulated = cross-family by defn),
%     preserves tier resolution, adoptable. RECOMMENDED.
%   strict cap @ Causally Suggestive -> louder normative statement but reviewers (esp. MI
%     empiricists, R1) read it as rigged to produce low verdicts; collapses tier resolution
%     (everything piles at Causally Suggestive). NOT recommended.
% OPEN QUESTION FOR JENSEN: can preregistration ever SUBSTITUTE for a missing evidence family,
% or only BOOST within-family severity? (design call for a causal-methodology expert.)
\TODO{Write up ~1 paragraph after Jensen weighs in. Keep verdict LABELS clean (no provenance
suffix); discuss corroboration/preregistration status in the per-claim PROSE and in this
subsection. Keep a 2-sentence confirmation/corroboration grounding here (full version in supplement).}

\subsection{Ablation and Patching Are Different Interventions}
% R3: make one piece of the causal apparatus actually load-bearing.
\PORT{The do-operator distinction from v5 \S7.3 (ablation = do(h:=0/E[h]); patching =
do(h:=h')) and use it to REINTERPRET the method-conditional faithfulness finding (E1).}
\TODO{This is the only formal causal content that stays in the main paper. Make it derive
something the prose didn't already say. Move the decorative IOI SCM to the supplement.}

% =====================================================================
\section{Audit: Sixteen Mechanistic Claims}
\label{sec:case-studies}
% THE PAYOFF. Reframed as claim-level.
\TODO{Open with the framing paragraph: "We audit CLAIMS, not papers. A verdict characterizes
the current evidential status of a claim, aggregating all published evidence to date, and is
not an assessment of any single paper's quality. Verdicts are graded against an idealized
standard; a low tier marks distance from that ideal, not a flawed study." (fixes R1)}
\PORT{tab:case-studies (13-circuit verdict table) + the 3--4 illustrative case studies
(Induction Heads=Triangulated, IOI=Mechanistically Supported, Probing=Proposed,
Gender Bias=Disconfirmed) + the "Recurring Validity Gaps" subsection.}
\TODO{For each illustrative case, show a short SCORING TRACE (which criteria pass/partial/
untested and why) so verdicts are visibly falsifiable, not asserted (R1). Cite follow-up
evidence per claim (this is the per-claim mini-lit-review) --- see lit_review_findings.md.}
\TODO{IOI enrichment: Miller/Chughtai 2024 (faithfulness not robust, node-vs-edge -> E1), ACDC +
"Demystifying Variance" (non-uniqueness), Adaptive Circuit Behavior (E2), Adversarial Circuit Eval
ICLR 2024, Adhikari 2026 (minimal circuits -> non-uniqueness/minimality).
CROSS-SCALE/CROSS-TASK (new E-evidence): Merullo 2024 (ALREADY in bib as a synthesis method) shows
IOI reproduces on a LARGER GPT-2 (E4 untested -> PARTIAL) and ~78% head reuse on Colored Objects
(E3) --- nudges the IOI trace but likely stays Mechanistically Supported.
SELF-REPAIR caveat (pushes the OTHER way): Conditional Co-Ablation 2026 + Self-Repair survey 2026
= MI-NATIVE evidence that backup name-mover claims can be ablation-protocol ARTIFACTS. Grounds v5's
EXISTING system-reserve/backup theme (pharma lens, "Backup circuits compensate after ablation") in
MI evidence, not just analogy; reinforces E1. Consider a short field-wide-ablation-confound paragraph.
Induction heads: "Induction Head Implementation Across Diverse Architectures" = E4 corroboration.
Probing: "Dissociating Decodability and Causal Use" = cleanest "decodable != used" cite.}
\TODO{GENDER BIAS = Disconfirmed --- handle carefully (thinnest evidence, most provocative verdict).
No 2024-2026 work re-tests separability of "bias" from grammatical-gender knowledge; 5 years on it
is STILL framed as open (2025 "Universal Patterns of Grammatical Gender"; 2026 disentanglement
preprint). Reframe: Disconfirmed on C4 --- the construct PRESUPPOSES a separation Vig et al.'s own
mediation contradicts (shared heads) and NO work has established. Do NOT claim to have proven
non-separability. OPTION: make the separability test a NEW Sec 6 experiment (thinnest verdict ->
empirical; shows the framework generating a novel test). Decide with Jensen.}

\subsection{Non-Circuit Claims}
\label{sec:noncircuit}
% R9/R2 method-agnosticism: two FULL non-circuit audits (Elliot's decision: do both, not a
% comparative paragraph). Together they show the framework produces DIFFERENT structured
% profiles for different method types --- the strong generality claim: observational SAE
% features cap on the CAUSAL axis; interventional steering vectors clear the causal bar but
% cap on the INTERPRETIVE (mode-match) axis.

\subsubsection{SAE Feature}
% Evidence base found (lit_review_findings.md). Pivot to the CANONICAL-UNITS claim.
\TODO{Audit "SAE features are canonical, monosemantic units corresponding to concepts." Evidence:
(1) NOT atomic / NOT canonical -> C4 discriminant failure: Leask et al. ICLR 2025 (stitching /
meta-SAEs) --- an "Einstein" latent decomposes into "scientist"+"Germany"+"famous person" (the
centerpiece example);
(2) NOT reliable/stable -> M1/M3: latent identity varies by training SEED (seed-dependence paper);
(3) decodable != used -> Braun et al. 2024 (end-to-end dict learning targets FUNCTIONALLY
important features, conceding standard training doesn't guarantee use);
(4) MIB (ICML 2026): SAE features perform NO BETTER than raw neurons on causal variable
localization (supervised DAS wins) --- direct causal-use ammunition.
Predicted verdict: Proposed / Causally Suggestive, capped by C4 + M1/M3.}

\subsubsection{Steering Vector}
% Evidence base found (lit_review_findings.md). Target: the refusal direction.
\TODO{Audit "adding the refusal direction v steers refusal" (Arditi et al. 2024, "Refusal is
mediated by a single direction"). Evidence:
(1) interventional (do(h += v)) -> clears sufficiency-to-steer that probing/SAEs cannot; strong
necessity+sufficiency in the original;
(2) dose-response (E5) by scaling v --- but "Inverse Scaling in Activation Steering" shows
effectiveness degrades with scale (E5 boundary condition);
(3) mode-match FAILURE (V2), confirmed by the literature's OWN drift: "Beyond a Single Direction:
CoT Disrupts Simple Steering of Refusal" (2026), adversarial-robustness + "Predicting Where
Steering Vectors Succeed" (E5/E6 fragility), COSMIC 2025 (original not fully general -> E4). The
field increasingly narrates refusal as a behavioral LEVER, not a representation --- exactly V2.
Predicted verdict: Mechanistically Supported (behavioral axis), capped by V2 + E5/E6 fragility.}

\subsubsection{J-space / "Global Workspace" (Anthropic 2026) --- flagship mode-conflation case}
% Full detail: jspace_case_notes.md. MechVal SLICE only here; cross-framework (invariance-depth,
% unification, MechViews) is routed to the OTHER papers.
\TODO{Audit Anthropic's J-space claim (verbalizable "global workspace" subspace via the Jacobian
lens; transformer-circuits.pub/2026/workspace). Predicted verdict: MECHANISTICALLY SUPPORTED (same
tier as IOI) --- necessity + sufficiency via three converging interventions (swap/ablate/inject);
caps below Triangulated on (a) single evidence family [C3: Nanda's Qwen replication adds a MODEL,
not a FAMILY], (b) I3 specificity UNTESTED (does ablation hurt only deliberate tasks or every
high-coordination task? = cardiac-stent problem), (c) I4 untested, (d) V4 interpretive inflation.
THE flagship description-mode conflation (implementational-statistical -> representational ->
computational "workspace") --- better than Othello: timely, high-profile, safety-relevant. Sharp
extra point: the measurement OPERATIONALIZES the construct (J-space = defined-as-reportable), so the
deliberate-reasoning correlation is guaranteed by construction (C4/C1 circularity). C5 weak: GWT
contested (COGITATE 2025 GNW prediction failed).
TWO-TIER SPLIT (one paper, multiple claims --- the showcase for "audit claims, not papers"): the
PARENT "global workspace" claim = Mechanistically Supported (capped by V4 + I3). The NARROWER claim
"the J-space carries deception-relevant signal detectable before output" sits CLOSER TO TRIANGULATED
--- explicit contrastive foil (aligned vs deceptive internal state) = more direct interventional
logic, AND it needs NO GWT narrative, so it dodges the V4 inflation and untested specificity that cap
the parent. Feed the deception SUB-claim into Sec 5 (safety monitor / minimum deployment; Nanda: needs
replication before trusted). CONSIDER promoting J-space to the running example in the Description-Modes
section (3.1) and/or the V4 flagship alongside/instead of Othello.}

\TODO{Update all "thirteen"/"13"/"fifteen"/"15" counts to SIXTEEN/16 (abstract, section title,
tab:case-studies caption, recurring-gaps prose).}

% =====================================================================
\section{Why It Matters}
\label{sec:implications}
% The "so what". Keep tight. Mark safety claims explicitly as PROPOSALS (R5).
\PORT{Safety applications from v5 discussion: chain-of-thought faithfulness (4 claims not 1),
deception detection, safety-evidence-gaps-at-scale, minimum deployment standards.}
\TODO{Frame deployment tiers as PROPOSALS derived from the dependency structure, not empirical
results (R5). If possible, one worked example where the tier predicts a monitor's failure.}
\PORT{Adoption mechanisms (validity statements, reviewer checklists, evidence cards) and the
five concrete recommendations (test specificity by default, etc.).}

\subsection{Standards for the Field}
% Elliot's decision: make the paper PRESCRIPTIVE (how mechanistic science SHOULD be conducted),
% not only diagnostic. GUARDRAIL: derive every "should" from (a) the dependency chain (logical
% necessity --- you cannot trust external validity without internal validity) and (b) PRECEDENT
% in sciences that already crossed this bridge (CONSORT for trials, registered reports in
% psychology, genome-wide significance + replication in GWAS). NEVER prescribe from authority/
% taste --- a solo unknown prescribing to the field triggers status-defense (R1/R2). Derived
% prescription reads as "the field's own standards made explicit," which is earned.
\TODO{Frame MI as where neuroscience/genetics/medicine were BEFORE they adopted these standards;
propose preregistration + corroboration + cross-family evidence as the concrete standard, each
traced to logic or precedent. Positioning shift: diagnostic tool -> call for standards (stronger,
more citable nugget; fits the org identity). Confirm the shift is what you want before finalizing.}
\TODO{Surface the practitioner 3-question path prominently here (R8): Is the description mode
declared? Is the claimed tier supported? Does evidence tier match claim tier?}

% =====================================================================
\section{Validating the Framework Itself}
\label{sec:self-validation}
% NEW SECTION. This is what flips "plausible rubric" -> "tested instrument" and is the
% single most important addition (R4, R1, R8, R10). Solo-authored, preregistered.
% CONFIRMED (Elliot, decisions 1 & 2): do BOTH the ground-truth Validated case AND the
% preregistered prospective audit --- not a fallback protocol. Advisor may give feedback
% (feedback only; authorship stays solo).

\subsection{A Ground-Truth Validated Claim}
% Elliot ran the C8 search. Clock/Pizza beats Tracr as a calibration design (see lit_review_findings.md).
\TODO{PRIMARY calibration case: grokking modular addition with the Clock vs Pizza dual mechanism
(Nanda 2023; "Clock and the Pizza," NeurIPS 2023). Two DISTINCT fully-known ground-truth
mechanisms + a Fourier diagnostic (mode at 2k mod p = Pizza, k = Clock) let us test not just
"can a claim reach Validated" but whether the tiers DISCRIMINATE: claim matches ground truth A
vs B vs neither. Produce the first Validated-tier verdict AND show a wrong-mechanism claim is
correctly denied it. SECONDARY: Tracr (Lindner 2023, RASP->weights) + Neural Decompiling of Tracr
(2024). (fixes R4 "nothing executed")}

\subsection{Preregistered Prospective Audit}
\TODO{Preregister verdicts on N newly published circuits BEFORE reviewing their evidence
(commit SHA, frozen). Report how the prospective verdicts land. This removes the retrofitting
circularity you already flagged in v5 limitations (R1/R4) and is your signature discipline.}

\subsection{Inter-Rater Reliability}
\TODO{Report agreement between >=2 independent raters on the 13 case-study verdicts (and/or
criterion-level agreement). Emphasize the structured profile --- not the tier --- as the
primary output, so disagreement localizes to individual falsifiable criteria (R8).}

\TODO{Optional: operationalize preregistration as a SEVERITY modulator (Mayo) inside criterion
status --- a preregistered specificity test is more severe than a post-hoc one, so it can raise
a criterion's status. This is where preregistration earns a place INSIDE the framework (R7).}

% =====================================================================
\section{Scope, Limitations, and Conclusion}
\label{sec:limitations}
\PORT{v5 limitations KEEPERS: inter-rater reliability, "case studies are illustrative not
preregistered" (now partly resolved by Sec 6), tier thresholds provisional, falsifiability
of the framework itself.}
\TODO{Scope statement (R9/R2): demonstrated claims now span attention-head circuits, an SAE
feature, and a steering vector (Sec 4.x) --- so method-agnosticism is SHOWN, not conjectured.
Transcoders / crosscoders remain conjecture with a stated adaptation path (supplement).}
\TODO{Beyond-MI generality: use a DESIGN-generality sentence (no cite, no commitment, cannot be
challenged) --- "Nothing in these criteria is specific to interpretability; 'is the construct
well-defined?' and 'is the evidence causal?' are domain-general, so the framework extends
naturally to mechanistic claims in other empirical sciences." Do NOT say "in prep" or cite the
chemistry audit (deanonymization + invites "can we see it?"). Save chemistry for the Perspective.}
\PORT{Conclusion paragraph.}

% =====================================================================
\bibliographystyle{tmlr}
\nocite{campbell1963experimental,loevinger1957objective}
\bibliography{references}

\end{document}
